% ============================================================
% TUTORIAL: Fase 1 --- Backend Core con AWS SAM + Lambda + DynamoDB
% Proyecto: Sentinel AcademIA
% Compilar: tectonic tutorial-fase1-aws-sam.tex
% ============================================================

\documentclass[11pt,a4paper,oneside]{article}

\usepackage[utf8]{inputenc}
\usepackage[T1]{fontenc}
\usepackage[spanish,es-nodecimaldot]{babel}
\usepackage{lmodern}
\usepackage[margin=2.5cm]{geometry}
\usepackage{parskip}
\usepackage{microtype}

\usepackage{xcolor}
\usepackage{colortbl}
\usepackage{array}
\usepackage{booktabs}
\usepackage{amssymb}
\usepackage{pifont}
\usepackage{tcolorbox}
\tcbuselibrary{breakable,skins}

\usepackage{listings}

\lstdefinelanguage{yaml}{
  basicstyle=\ttfamily\small,
  keywordstyle=\color{blue!70!black}\bfseries,
  stringstyle=\color{red!70!black},
  commentstyle=\color{green!50!black}\itshape,
  morekeywords={AWSTemplateFormatVersion,Transform,Description,Parameters,Globals,Resources,Outputs,Type,Properties,Environment,Variables,Events},
  literate=
    *{\{}{{{\{}}}{1}
     {\}}{{{\}}}}{1}
     {[}{{{[}}}{1}
     {]}{{{]}}}{1}
}

\lstdefinelanguage{bash}{
  basicstyle=\ttfamily\small,
  keywordstyle=\color{blue!70!black}\bfseries,
  stringstyle=\color{red!70!black},
  commentstyle=\color{green!50!black}\itshape,
  morekeywords={aws,sam,cd,ls,mkdir,rm,cp,mv,export,unset,echo,cat,grep},
}

\lstdefinelanguage{json}{
  basicstyle=\ttfamily\small,
  showstringspaces=false,
  breaklines=true,
  literate=
    *{\{}{{{\{}}}{1}
     {\}}{{{\}}}}{1}
     {[}{{{[}}}{1}
     {]}{{{]}}}{1}
     {:}{{{:}}}{1}
     {,}{{{,}}}{1},
  string=[s]{"}{"},
  stringstyle=\color{red!70!black},
  morestring=[b]",
}

\lstset{
  basicstyle=\ttfamily\small,
  numberstyle=\tiny\color{gray},
  numbers=left,
  numbersep=8pt,
  frame=single,
  framerule=0.4pt,
  rulecolor=\color{gray!50},
  backgroundcolor=\color{gray!5},
  breaklines=true,
  breakatwhitespace=true,
  showstringspaces=false,
  captionpos=b,
  tabsize=2,
  literate={á}{{\'a}}1 {é}{{\'e}}1 {í}{{\'i}}1 {ó}{{\'o}}1 {ú}{{\'u}}1
           {ñ}{{\~n}}1 {Á}{{\'A}}1 {É}{{\'E}}1 {Í}{{\'I}}1 {Ó}{{\'O}}1 {Ú}{{\'U}}1
           {¿}{{?`}}1 {¡}{{!`}}1
}

\usepackage[colorlinks=true,linkcolor=blue!60!black,urlcolor=blue!60!black,citecolor=blue!60!black]{hyperref}

\usepackage{fancyhdr}
\pagestyle{fancy}
\fancyhf{}
\fancyhead[L]{\small\textsc{Sentinel AcademIA}}
\fancyhead[R]{\small Tutorial Fase 1}
\fancyfoot[C]{\thepage}
\renewcommand{\headrulewidth}{0.4pt}

\definecolor{okgreen}{RGB}{40,140,60}
\definecolor{errred}{RGB}{200,50,50}
\definecolor{warnorange}{RGB}{220,130,30}
\definecolor{infoblue}{RGB}{30,90,180}
\definecolor{codebg}{RGB}{245,245,245}

\newcommand{\cmd}[1]{\texttt{\small #1}}
\newcommand{\filename}[1]{\texttt{\small #1}}
\newcommand{\ok}{\textcolor{okgreen}{\textbf{\checkmark}}}
\newcommand{\err}{\textcolor{errred}{\textbf{\ding{55}}}}
\newcommand{\warn}{\textcolor{warnorange}{\textbf{\textasciitilde}}}

\title{\Huge\textbf{Fase 1: Backend Core} \\[0.3em]
       \Large AWS SAM + Lambda + DynamoDB Multi-Tenant \\[0.5em]
       \normalsize Tutorial paso a paso}
\author{Proyecto Sentinel AcademIA \\ \small lFunknown}
\date{\today}

\begin{document}

\maketitle
\thispagestyle{empty}

\begin{tcolorbox}[colback=blue!5!white,colframe=infoblue,title={\textbf{¿Qué construimos en este tutorial?}}]
Este documento es el log completo de la \textbf{Fase 1} del proyecto Sentinel AcademIA: el backend serverless en AWS Academy con Lambdas Python, DynamoDB multi-tenant, y API Gateway.

\textbf{Stack tecnol\'ogico:}
\begin{itemize}
  \item Python 3.12 (runtime de Lambda) + 3.14.3 (local)
  \item AWS SAM (Infrastructure as Code)
  \item Lambda + API Gateway REST + DynamoDB
  \item Pydantic v2 para validaci\'on
  \item aws-lambda-powertools para observabilidad
  \item OpenAPI spec como contrato
\end{itemize}

\textbf{Lo que ver\'as:}
\begin{itemize}
  \item Estructura de carpetas y archivos
  \item Cada m\'odulo compartido con su prop\'osito
  \item El template SAM con todas las decisiones explicadas
  \item Las trampas de AWS Academy (y c\'omo resolverlas)
  \item 5 smoke tests que validan todo end-to-end
\end{itemize}
\end{tcolorbox}

\begin{tcolorbox}[colback=warnorange!5!white,colframe=warnorange,title={\textbf{Aviso: AWS Academy es restrictivo}}]
Este proyecto corre en \textbf{AWS Academy Learner Lab}, que tiene l\'imites que no aplican a AWS normal. En particular, \textbf{no podemos crear IAM roles}. Esto afecta TODO el flujo y se explica en detalle en la secci\'on correspondiente.
\end{tcolorbox}

\tableofcontents
\newpage

% ============================================================
\section{Prerrequisitos}
% ============================================================

\subsection{Herramientas instaladas}

\begin{center}
\begin{tabular}{@{}lll@{}}
\toprule
\textbf{Herramienta} & \textbf{Versi\'on} & \textbf{Check} \\
\midrule
AWS CLI & 2.35+ & \cmd{aws --version} \\
SAM CLI & 1.162+ & \cmd{sam --version} \\
Python 3.12 & 3.12.13 & \cmd{python3.12 --version} \\
Python 3.14.3 & 3.14.3 & \cmd{python3 --version} \\
mise & 2026.6+ & \cmd{mise --version} \\
Docker & 29+ & \cmd{docker --version} \\
\bottomrule
\end{tabular}
\end{center}

\subsubsection{Instala Python 3.12 con mise (si no lo tienes)}

AWS Lambda runtime es Python 3.12. Tu sistema puede tener 3.14 (v\'ia OCI installer). Necesitas 3.12 tambi\'en:

\begin{lstlisting}[language=bash]
mise use --global python@3.12
mise exec -- python3.12 --version  # debe decir 3.12.13
\end{lstlisting}

\subsection{Credenciales de AWS Academy}

Las credenciales \textbf{no} se guardan en archivos del proyecto. Viven solo en variables de entorno del shell:

\begin{lstlisting}[language=bash]
export AWS_ACCESS_KEY_ID="ASIAT..."
export AWS_SECRET_ACCESS_KEY="..."
export AWS_SESSION_TOKEN="IQoJb3JpZ2luX2Vj..."
export AWS_DEFAULT_REGION="us-east-1"
export AWS_REGION="us-east-1"

aws sts get-caller-identity
# debe mostrar: Account "227165337884" o el de tu lab
\end{lstlisting}

\textbf{Vida \'util}: 3-4 horas. Cuando expiren, repites el flujo.

% ============================================================
\section{Estructura del proyecto}
% ============================================================

\begin{lstlisting}[language=bash]
sentinel-academia/
|-- infra/
|   |-- template.yaml           # SAM template (IaC)
|   `-- samconfig.toml          # Configuracion de sam deploy
|
|-- src/
|   |-- handlers/               # 1 archivo = 1 Lambda
|   |   |-- _shared/            # Modulos compartidos entre Lambdas
|   |   |   |-- __init__.py     # <-- IMPORTANTE: marca como package
|   |   |   |-- config.py
|   |   |   |-- logger.py
|   |   |   |-- errors.py
|   |   |   |-- http_response.py
|   |   |   |-- correlation_id.py
|   |   |   |-- dynamo_client.py
|   |   |   |-- schemas.py
|   |   |   `-- auth.py
|   |   |-- __init__.py         # <-- IMPORTANTE: marca como package
|   |   |-- health_check.py
|   |   |-- create_queja.py
|   |   |-- get_queja.py
|   |   `-- list_quejas.py
|   |
|   `-- requirements.txt        # Runtime deps (las usa SAM)
|
|-- pyproject.toml              # Configuracion completa (dev + runtime)
|-- .opencode/
|   `-- rules/10-aws-academy.md # Regla especifica de Academy
`-- api-mock/openapi.yaml       # Spec de la API
\end{lstlisting}

\subsection{Decisiones clave de la estructura}

\begin{itemize}
  \item \ok \textbf{1 archivo = 1 Lambda}: cada endpoint vive en su propio archivo. F\'acil de testear, deployar, y razonar.
  \item \ok \textbf{\_shared/ para c\'odigo com\'un}: 8 m\'odulos que usan TODAS las Lambdas.
  \item \ok \textbf{infra/ separado de src/}: el template vive aparte del c\'odigo, evita ciclos de importaci\'on.
  \item \ok \textbf{\_\_init\_\_.py en cada package}: sin esto, Python no reconoce la carpeta como package.
  \item \err \textbf{NO usamos src/handlers/ como prefijo}: los handlers son top-level para que Lambda los importe f\'acil.
\end{itemize}

% ============================================================
\section{Configuraci\'on inicial: pyproject.toml}
% ============================================================

El archivo \filename{pyproject.toml} centraliza la configuraci\'on del proyecto Python: dependencias, herramientas de calidad, y versionado.

\begin{lstlisting}[language=Python,caption={pyproject.toml (deps runtime)}]
[project]
name = "sentinel-academia"
version = "0.1.0"
requires-python = ">=3.12,<3.15"  # Lambda = 3.12, local = 3.14
dependencies = [
    "aws-lambda-powertools[tracer,logger,metrics]>=3.0.0",
    "boto3>=1.40.0",
    "boto3-stubs[dynamodb,s3,ses,sns,sqs]>=1.40.0",
    "pydantic>=2.8.0",
    "pydantic-settings>=2.4.0",
    "orjson>=3.10.0",
]

[project.optional-dependencies]
dev = [
    "pytest>=8.3.0",
    "moto[dynamodb,s3,ses,sns,sqs,secretsmanager]>=5.0.0",
    "ruff>=0.6.0",
    "mypy>=1.11.0",
]
\end{lstlisting}

\subsection{requirements.txt para SAM}

AWS SAM \textbf{NO lee pyproject.toml} para dependencias. Necesita un \filename{requirements.txt} plano:

\begin{lstlisting}[language=bash,caption={src/requirements.txt}]
# Runtime dependencies for Sentinel AcademIA Lambdas
aws-lambda-powertools[tracer,logger,metrics]>=3.0.0
boto3>=1.40.0
boto3-stubs[dynamodb]>=1.40.0
pydantic>=2.8.0
pydantic-settings>=2.4.0
orjson>=3.10.0
email-validator>=2.2.0  # <-- IMPORTANTE para EmailStr de Pydantic
\end{lstlisting}

\begin{tcolorbox}[colback=errred!5!white,colframe=errred,title={Trampa \#1: email-validator}]
Si usas \texttt{EmailStr} de Pydantic, necesitas \texttt{email-validator}. Si no, en runtime:
\begin{lstlisting}
Runtime.ImportModuleError: email-validator is not installed,
  run `pip install 'pydantic[email]'`
\end{lstlisting}
\end{tcolorbox}

% ============================================================
\section{Los 8 m\'odulos compartidos (\_shared/)}
% ============================================================

\subsection{config.py --- Configuraci\'on desde env vars}

\begin{lstlisting}[language=Python,caption={handlers/\_shared/config.py (resumen)}]
from functools import lru_cache
from pydantic_settings import BaseSettings

class Settings(BaseSettings):
    app_name: str = "sentinel-academia"
    stage: str = "dev"
    region: str = "us-east-1"
    dynamodb_table: str = "sentinel-quejas-{stage}"
    files_bucket: str = "sentinel-files-{stage}"
    log_level: str = "INFO"
    allowed_origins: str = "*"
    # ... OCI vars (Fase 2)

@lru_cache(maxsize=1)
def get_settings() -> Settings:
    return Settings()  # lee de env vars automaticamente
\end{lstlisting}

\textbf{Conceptos clave:}
\begin{itemize}
  \item \ok Pydantic \texttt{BaseSettings} lee env vars \textbf{autom\'aticamente}
  \item \ok \texttt{@lru\_cache} garantiza un singleton por invocaci\'on Lambda
  \item \ok \texttt{stage} permite tener dev/staging/prod en la misma cuenta
  \item \err Nunca hardcodees valores; siempre desde env vars
\end{itemize}

\subsection{logger.py --- Logger estructurado con powertools}

\begin{lstlisting}[language=Python,caption={handlers/\_shared/logger.py}]
from aws_lambda_powertools import Logger

_loggers: dict[str, Logger] = {}

def get_logger(name: str) -> Logger:
    if name in _loggers:
        return _loggers[name]
    settings = get_settings()
    logger = Logger(
        service=settings.app_name,
        level=settings.log_level,
    )
    _loggers[name] = logger
    return logger
\end{lstlisting}

\textbf{Reglas:}
\begin{itemize}
  \item \err \textbf{NUNCA uses print()} --- siempre \texttt{logger.info()}
  \item \ok El logger incluye autom\'aticamente correlation\_id y tenant\_id
  \item \ok Niveles: DEBUG, INFO, WARNING, ERROR
\end{itemize}

\subsection{errors.py --- Excepciones tipadas + error handler}

\begin{lstlisting}[language=Python,caption={handlers/\_shared/errors.py (nucleo)}]
class AppError(Exception):
    def __init__(self, message, status_code=500, error_code="INTERNAL_ERROR",
                 details=None):
        super().__init__(message)
        self.message = message
        self.status_code = status_code
        self.error_code = error_code
        self.details = details or {}

class ValidationError(AppError):  # 400
    def __init__(self, message, details=None):
        super().__init__(message, 400, "VALIDATION_ERROR", details)

class UnauthorizedError(AppError):  # 401
    def __init__(self, message="Unauthorized"):
        super().__init__(message, 401, "UNAUTHORIZED")

class NotFoundError(AppError):  # 404
    ...

def error_handler(func):
    def wrapper(event, context):
        try:
            return func(event, context)
        except AppError as e:
            log.warning("AppError raised", extra={
                "error_code": e.error_code,
                "status_code": e.status_code,
                "error_message": e.message,  # <-- no "message"!
            })
            return build_error_response(
                status_code=e.status_code,
                error_code=e.error_code,
                message=e.message,
                details=e.details,
            )
        except Exception as e:
            log.exception("Unhandled exception")
            return build_error_response(500, "INTERNAL_ERROR", ...)
    return wrapper
\end{lstlisting}

\begin{tcolorbox}[colback=errred!5!white,colframe=errred,title={Trampa \#2: KeyError 'message' en logger}]
El m\'odulo \texttt{logging} de Python tiene un atributo interno \texttt{message} en LogRecord. Si pasas \texttt{extra=\{``message'': ...\}} al logger, falla con \texttt{KeyError: ``Attempt to overwrite 'message' in LogRecord''}.

\textbf{Soluci\'on}: usa \texttt{error\_message} o cualquier otro nombre.
\end{tcolorbox}

\subsection{http\_response.py --- Respuestas estandar}

\begin{lstlisting}[language=Python,caption={handlers/\_shared/http\_response.py (fragmento)}]
class HttpResponse:
    @staticmethod
    def ok(body): return build_response(200, body)
    @staticmethod
    def created(body): return build_response(201, body)
    @staticmethod
    def accepted(body): return build_response(202, body)  # <-- async pattern
    @staticmethod
    def bad_request(msg, details=None):
        return build_error_response(400, "VALIDATION_ERROR", msg, details)
    # ... not_found, conflict, internal_error, etc.

def build_error_response(status_code, error_code, message, details=None):
    body = {
        "code": error_code,
        "message": message,
        "correlationId": get_correlation_id() or "",
        "timestamp": datetime.now(timezone.utc).isoformat(),
    }
    if details:
        body["details"] = details
    return build_response(status_code, body)
\end{lstlisting}

\textbf{Decisiones:}
\begin{itemize}
  \item \ok \textbf{CORS siempre}: headers en TODAS las respuestas
  \item \ok \textbf{orjson} para serializar (m\'as r\'apido que stdlib)
  \item \ok Formato de error \textbf{alineado con OpenAPI spec}: \texttt{\{code, message, correlationId, timestamp\}}
\end{itemize}

\subsection{correlation\_id.py --- Trazabilidad end-to-end}

\begin{lstlisting}[language=Python,caption={handlers/\_shared/correlation\_id.py}]
import uuid
from contextvars import ContextVar

_correlation_id_var: ContextVar[str | None] = ContextVar(
    "correlation_id", default=None
)

def with_correlation_id(event):
    headers = event.get("headers", {}) or {}
    cid = (headers.get("x-correlation-id")
           or headers.get("X-Correlation-ID")
           or headers.get("X-Correlation-Id"))
    if not cid:
        cid = str(uuid.uuid4())
    _correlation_id_var.set(cid)
    return cid

def get_correlation_id() -> str | None:
    return _correlation_id_var.get()
\end{lstlisting}

\textbf{Conceptos:}
\begin{itemize}
  \item \ok \texttt{ContextVar}: variable local al contexto de ejecuci\'on (no global)
  \item \ok Header \texttt{X-Correlation-ID} opcional; se genera UUID si falta
  \item \ok Aparece en logs y en respuestas (header \texttt{X-Correlation-ID})
\end{itemize}

\subsection{dynamo\_client.py --- Wrapper multi-tenant}

\begin{lstlisting}[language=Python,caption={handlers/\_shared/dynamo\_client.py (key builders)}]
def tenant_pk(tenant_id: str) -> str:
    return f"TENANT#{tenant_id}"

def queja_pk(tenant_id: str, queja_id: str) -> str:
    return f"{tenant_pk(tenant_id)}#QUEJA#{queja_id}"

def queja_sk(queja_id: str) -> str:
    return f"QUEJA#{queja_id}"

class DynamoClient:
    def put_item(self, item, condition_expression=None):
        # PK/SK ya vienen en el item
        kwargs = {"Item": item}
        if condition_expression:
            kwargs["ConditionExpression"] = condition_expression
        self._table.put_item(**kwargs)

    def get_queja(self, tenant_id, queja_id):
        item = self.get_item(
            pk=queja_pk(tenant_id, queja_id),
            sk=queja_sk(queja_id)
        )
        if item is None:
            raise NotFoundError(...)
        return item

    def query_by_tenant(self, tenant_id, sk_prefix=None, limit=None):
        kc = Key("pk").eq(tenant_pk(tenant_id))
        if sk_prefix:
            kc = kc & Key("sk").begins_with(sk_prefix)
        return self._table.query(KeyConditionExpression=kc, ...)["Items"]
\end{lstlisting}

\textbf{Patrones clave:}
\begin{itemize}
  \item \ok \textbf{Single-table}: PK = \texttt{TENANT\#xxx\#QUEJA\#yyy}, SK = \texttt{QUEJA\#yyy}
  \item \ok \textbf{Multi-tenant safety}: las queries SIEMPRE filtran por tenant
  \item \ok \textbf{GSI para queries globales}: gsi1pk = \texttt{TENANT\#xxx\#STATUS\#PENDIENTE}
\end{itemize}

\subsection{schemas.py --- Validaci\'on con Pydantic}

\begin{lstlisting}[language=Python,caption={handlers/\_shared/schemas.py (modelos principales)}]
from enum import StrEnum

class QuejaStatus(StrEnum):
    PENDIENTE = "PENDIENTE"
    EN_COLA = "EN_COLA"
    PROCESANDO = "PROCESANDO"
    ANALIZADA = "ANALIZADA"
    NOTIFICADA = "NOTIFICADA"
    ERROR = "ERROR"

class CategoriaEnum(StrEnum):
    ACADEMICA = "ACADEMICA"
    INFRAESTRUCTURA = "INFRAESTRUCTURA"
    ACOSO = "ACOSO"
    ADMINISTRATIVA = "ADMINISTRATIVA"
    SALUD = "SALUD"
    OTRA = "OTRA"

class QuejaCreate(QuejaBase):
    titulo: str = Field(min_length=5, max_length=120)
    descripcion: str = Field(min_length=20, max_length=5000)
    categoriaDeclarada: CategoriaEnum = CategoriaEnum.OTRA
    anonima: bool = False
    sede: str | None = None
    facultad: str | None = None
    contactoEmail: EmailStr | None = None
    contactoTelefono: str | None = Field(default=None, pattern=r"^\+?[0-9\s\-]{7,20}$")

class QuejaAccepted(BaseModel):
    quejaId: str
    status: QuejaStatus = QuejaStatus.PENDIENTE
    correlationId: str
    createdAt: str
    estimatedAnalysisTime: int = 30
\end{lstlisting}

\textbf{Decisiones:}
\begin{itemize}
  \item \ok \textbf{Nombres de campos en camelCase} (alineados con OpenAPI)
  \item \ok \textbf{Enums UPPERCASE} (alineados con spec)
  \item \ok \textbf{Validaciones inline}: min\_length, max\_length, pattern, etc.
  \item \ok \texttt{QuejaAccepted} separado de \texttt{QuejaResponse} (uno para 202, otro para 200)
\end{itemize}

\subsection{auth.py --- Multi-tenant}

\begin{lstlisting}[language=Python,caption={handlers/\_shared/auth.py}]
def extract_tenant_id(event):
    headers = event.get("headers", {}) or {}
    tenant_id = (headers.get("X-Tenant-ID")
                 or headers.get("x-tenant-id")
                 or headers.get("X-Tenant-Id"))
    if not tenant_id:
        raise UnauthorizedError(
            "Missing X-Tenant-ID header. Multi-tenant access requires..."
        )
    tenant_id = tenant_id.strip()
    if not tenant_id or not all(c.isalnum() or c in "-_" for c in tenant_id):
        raise UnauthorizedError("Invalid X-Tenant-ID format")
    return tenant_id
\end{lstlisting}

\textbf{Conceptos:}
\begin{itemize}
  \item \ok Tenant ID viene de header \textbf{nunca} del body (seguridad)
  \item \ok Sanitizaci\'on: solo alfanum\'erico, gui\'on bajo, gui\'on
  \item \err Sin header $\rightarrow$ 401 UNAUTHORIZED
\end{itemize}

% ============================================================
\section{Las 4 Lambdas}
% ============================================================

\subsection{health\_check.py --- Endpoint sin auth}

\begin{lstlisting}[language=Python,caption={handlers/health\_check.py}]
from handlers._shared.errors import error_handler
from handlers._shared.http_response import HttpResponse
from handlers._shared.config import get_settings
from handlers._shared.correlation_id import with_correlation_id

@error_handler
def lambda_handler(event, context):
    with_correlation_id(event)
    settings = get_settings()
    return HttpResponse.ok({
        "status": "healthy",
        "service": settings.app_name,
        "stage": settings.stage,
        "region": settings.region,
    })
\end{lstlisting}

\subsection{create\_queja.py --- POST 202 Accepted (async pattern)}

\begin{lstlisting}[language=Python,caption={handlers/create\_queja.py (nucleo)}]
@error_handler
def lambda_handler(event, context):
    with_correlation_id(event)
    try:
        # 1. Auth: extraer tenant (multi-tenant safety)
        tenant_id = extract_tenant_id(event)
        user_id = extract_user_id(event)
        log.append_keys(tenant_id=tenant_id)
        if user_id:
            log.append_keys(user_id=user_id)

        # 2. Parse + validate body
        body = _parse_body(event)
        try:
            input_data = QuejaCreate.model_validate(body)
        except PydanticValidationError as e:
            errors = [{"field": "...", "message": "...", "type": "..."}
                      for err in e.errors()]
            raise ValidationError(
                "Request body validation failed",
                details={"errors": errors}
            ) from e

        # 3. Build accepted response
        now = now_iso()
        correlation_id = get_correlation_id() or ""
        accepted = QuejaAccepted(
            status="PENDIENTE",
            correlationId=correlation_id,
            createdAt=now,
            estimatedAnalysisTime=30,
        )

        # 4. Build DynamoDB item
        item = input_data.model_dump(mode="json")
        item["quejaId"] = accepted.quejaId
        item["tenantId"] = tenant_id
        item["userId"] = user_id if not input_data.anonima else None
        item["status"] = "PENDIENTE"
        item["createdAt"] = now
        item["updatedAt"] = now
        item["escalada"] = False
        item["pk"] = queja_pk(tenant_id, accepted.quejaId)
        item["sk"] = queja_sk(accepted.quejaId)
        item["gsi1pk"] = f"TENANT#{tenant_id}#STATUS#PENDIENTE"
        item["gsi1sk"] = now

        # 5. Persist (con condition: evita duplicados)
        dynamo = get_dynamo_client()
        dynamo.put_item(item, condition_expression="attribute_not_exists(pk)")

        # 6. Return 202 Accepted
        return HttpResponse.accepted(accepted.model_dump(mode="json"))
    finally:
        clear_correlation_id()
\end{lstlisting}

\textbf{Patrones importantes:}
\begin{itemize}
  \item \ok \textbf{202 Accepted} (no 201): procesamiento as\'incrono
  \item \ok \texttt{log.append\_keys(tenant\_id=...)}: a\~nade contexto a todos los logs
  \item \ok \texttt{condition\_expression="attribute\_not\_exists(pk)"}: previene duplicados
  \item \ok \texttt{try/finally} con \texttt{clear\_correlation\_id()}: limpia contexto
\end{itemize}

\subsection{get\_queja.py y list\_quejas.py}

Patrones similares:
\begin{itemize}
  \item Validan \texttt{pathParameters} (quejaId)
  \item Llaman a \texttt{dynamo.get\_queja()} o \texttt{query\_by\_tenant()}
  \item Retornan 200 con la queja o lista
  \item Manejan 404 con \texttt{NotFoundError}
\end{itemize}

% ============================================================
\section{SAM template --- Infraestructura como c\'odigo}
% ============================================================

\subsection{El template completo (infra/template.yaml)}

\begin{lstlisting}[language=yaml,caption={infra/template.yaml --- fragmentos clave}]
AWSTemplateFormatVersion: "2010-09-09"
Transform: AWS::Serverless-2016-10-31

Parameters:
  Stage:
    Type: String
    Default: dev
    AllowedValues: [dev, staging, prod]
  LambdaRoleArn:
    Type: String
    Default: "arn:aws:iam::227165337884:role/LabRole"

Globals:
  Function:
    Runtime: python3.12
    MemorySize: 256
    Timeout: 30
    Architectures: [x86_64]
    Tracing: Active
    Environment:
      Variables:
        STAGE: !Ref Stage
        LOG_LEVEL: INFO
        ALLOWED_ORIGINS: "*"
        POWERTOOLS_SERVICE_NAME: sentinel-academia
  Api:
    Cors:
      AllowMethods: "'GET,POST,PUT,DELETE,OPTIONS'"
      AllowHeaders: "'Content-Type,Authorization,X-Correlation-ID,X-Tenant-ID'"
      AllowOrigin: "'*'"

Resources:
  QuejasTable:
    Type: AWS::DynamoDB::Table
    Properties:
      TableName: !Sub "sentinel-quejas-${Stage}"
      BillingMode: PAY_PER_REQUEST
      AttributeDefinitions:
        - { AttributeName: pk,  AttributeType: S }
        - { AttributeName: sk,  AttributeType: S }
        - { AttributeName: gsi1pk, AttributeType: S }
        - { AttributeName: gsi1sk, AttributeType: S }
      KeySchema:
        - { AttributeName: pk, KeyType: HASH }
        - { AttributeName: sk, KeyType: RANGE }
      GlobalSecondaryIndexes:
        - IndexName: GSI1-StatusByDate
          KeySchema:
            - { AttributeName: gsi1pk, KeyType: HASH }
            - { AttributeName: gsi1sk, KeyType: RANGE }
          Projection: { ProjectionType: ALL }

  CreateQuejaFunction:
    Type: AWS::Serverless::Function
    Properties:
      Role: !Ref LambdaRoleArn           # <-- LabRole pre-existente
      FunctionName: !Sub "sentinel-create-queja-${Stage}"
      CodeUri: ../src                    # <-- desde la raiz del proyecto
      Handler: handlers.create_queja.lambda_handler
      MemorySize: 512
      Environment:
        Variables:
          DYNAMODB_TABLE: !Sub "sentinel-quejas-${Stage}"
      Events:
        CreateApi:
          Type: Api
          Properties:
            Path: /api/quejas           # <-- alineado con OpenAPI
            Method: POST
            RestApiId: !Ref SentinelApi
\end{lstlisting}

% ============================================================
\section{AWS Academy: las trampas y c\'omo resolverlas}
% ============================================================

\subsection{Regla de oro: SIEMPRE LabRole}

\begin{tcolorbox}[colback=warnorange!5!white,colframe=warnorange,title={\textbf{\ding{55} NO crees roles, \ding{51} usa LabRole}}]
En AWS Academy Learner Lab, el usuario \texttt{voclabs} \textbf{NO puede}:
\begin{itemize}
  \item \texttt{iam:CreateRole} --- denegado
  \item \texttt{iam:PassRole} --- solo sobre \texttt{LabRole}
  \item Tags en IAM resources --- denegado
\end{itemize}

\textbf{Consecuencia}: SAM normalmente auto-genera roles. Eso falla. Soluci\'on: declarar el role manualmente.
\end{tcolorbox}

\begin{lstlisting}[language=yaml,caption={C\'omo referenciar LabRole}]
LambdaRoleArn:
  Type: String
  Default: "arn:aws:iam::227165337884:role/LabRole"

CreateQuejaFunction:
  Properties:
    Role: !Ref LambdaRoleArn   # <-- en CADA Lambda
\end{lstlisting}

\subsection{No usar \texttt{Policies:} en las funciones}

Si pones:

\begin{lstlisting}[language=yaml]
Properties:
  Policies:
    - AWSLambdaBasicExecutionRole
\end{lstlisting}

SAM intenta crear un role nuevo $\rightarrow$ 403 AccessDenied.

\textbf{Soluci\'on}: no declarar Policies; el LabRole ya tiene los permisos necesarios.

\subsection{CodeUri: \texttt{../src} (no \texttt{../} ni \texttt{../../})}

El CodeUri determina qu\'e se incluye en el zip de la Lambda:

\begin{lstlisting}[language=yaml,caption={Comparaci\'on de CodeUri}]
# MAL: incluye TODO el proyecto (docs, web, AGENTS.md, .gitignore...)
CodeUri: ../

# MAL: incluye src/ + el template otra vez
CodeUri: ../src/../

# BIEN: solo src/, zip limpio
CodeUri: ../src
\end{lstlisting}

\textbf{Verificaci\'on}: tras \texttt{sam build}, el zip debe tener:
\begin{itemize}
  \item \filename{handlers/\_\_init\_\_.py} (marcando el package)
  \item \filename{handlers/\_shared/\_\_init\_\_.py}
  \item \filename{handlers/create\_queja.py}
  \item Site-packages (boto3, pydantic, etc.)
\end{itemize}

\subsection{El bug \texttt{No module named 'src'}}

Cuando el handler era \texttt{src.handlers.create\_queja.lambda\_handler}, Lambda dec\'ia:

\begin{lstlisting}
Runtime.ImportModuleError: No module named 'src'
\end{lstlisting}

\textbf{Causa}: aunque el zip ten\'ia \texttt{src/}, el handler busc\'o \texttt{src.handlers.*} y no encontr\'o \texttt{src} en el path de importaci\'on.

\textbf{Soluci\'on}: renombrar handlers para que sean \textbf{top-level}:

\begin{lstlisting}[language=bash]
# Cambiar todas las referencias en el codigo
sed -i 's|from src\.handlers|from handlers|g' src/handlers/**/*.py
sed -i 's|src\.handlers\.|handlers.|g' template.yaml
\end{lstlisting}

% ============================================================
\section{Build, deploy, y samconfig.toml}
% ============================================================

\subsection{samconfig.toml --- Configuraci\'on persistente}

\begin{lstlisting}[language=yaml,caption={infra/samconfig.toml}]
version = 0.1

[default.deploy.parameters]
stack_name = "sentinel-academia-dev"
s3_bucket = "sentinel-sam-artifacts-1781951709"
s3_prefix = "sentinel-academia-dev"
region = "us-east-1"
confirm_changes_before_deploy = false
capabilities = "CAPABILITY_IAM"
disable_rollback = false

[default.deploy.global.parameters]
debug = false
\end{lstlisting}

\subsection{El ciclo build $\rightarrow$ deploy}

\begin{lstlisting}[language=bash,caption={Comandos esenciales}]
# 1. Limpiar build anterior
rm -rf .aws-sam .aws-sam-new

# 2. Build (crea .aws-sam/build/<Function>/ con deps instaladas)
sam build --build-dir .aws-sam-new

# 3. Deploy (sube zips a S3 + crea/actualiza stack)
sam deploy --stack-name sentinel-academia-dev \
  --template-file .aws-sam-new/template.yaml \
  --s3-bucket sentinel-sam-artifacts-1781951709 \
  --capabilities CAPABILITY_IAM \
  --no-disable-rollback \
  --no-confirm-changeset

# 4. Ver outputs
aws cloudformation describe-stacks \
  --stack-name sentinel-academia-dev \
  --query 'Stacks[0].Outputs'
\end{lstlisting}

\begin{tcolorbox}[colback=errred!5!white,colframe=errred,title={Trampa \#3: \_\_pycache\_\_ con permisos de root}]
Tras \texttt{sam build} con Docker, se crean archivos \texttt{\_\_pycache\_\_/*.pyc} con owner \texttt{root}. No puedes borrarlos con tu usuario.

\textbf{Soluci\'on}: usa \texttt{-{}-build-dir} apuntando a un directorio nuevo:
\begin{lstlisting}
sam build --build-dir .aws-sam-new
\end{lstlisting}
\end{tcolorbox}

% ============================================================
\section{Los 5 smoke tests que validan todo}
% ============================================================

\begin{lstlisting}[language=bash,caption={tests/smoke.sh}]
API="https://om3eiczjr9.execute-api.us-east-1.amazonaws.com/dev"

# Test 1: Health check
curl -s "$API/health"
# Esperado: 200 + {"status":"healthy",...}

# Test 2: Crear queja (alineado con OpenAPI)
curl -s -X POST "$API/api/quejas" \
  -H "Content-Type: application/json" \
  -H "X-Tenant-ID: utpl" \
  -H "X-User-ID: estudiante-001" \
  -H "X-Correlation-ID: test-001" \
  -d '{"titulo":"Problema","descripcion":"...","categoriaDeclarada":"ADMINISTRATIVA","anonima":false}'
# Esperado: 202 + QuejaAccepted

# Test 3: Listar quejas
curl -s "$API/api/quejas" -H "X-Tenant-ID: utpl"
# Esperado: 200 + {items, total, nextCursor}

# Test 4: Obtener queja por ID
curl -s "$API/api/quejas/<id>" -H "X-Tenant-ID: utpl"
# Esperado: 200 + queja completa

# Test 5: Sin X-Tenant-ID (debe fallar con 401)
curl -s -X POST "$API/api/quejas" -H "Content-Type: application/json" -d '...'
# Esperado: 401 + {code:UNAUTHORIZED,message:...,correlationId:...,timestamp:...}
\end{lstlisting}

\subsection{Output esperado de cada test}

\begin{lstlisting}[language=json,caption={Test 1: GET /health}]
{
  "status": "healthy",
  "service": "sentinel-academia",
  "stage": "dev",
  "region": "us-east-1"
}
\end{lstlisting}

\begin{lstlisting}[language=json,caption={Test 2: POST /api/quejas}]
{
  "quejaId": "1f6733ea-0578-48cf-af10-740962f7ced0",
  "status": "PENDIENTE",
  "correlationId": "test-001",
  "createdAt": "2026-06-20T13:34:24.779106+00:00",
  "estimatedAnalysisTime": 30
}
\end{lstlisting}

\begin{lstlisting}[language=json,caption={Test 5: error 401 (formato spec-aligned)}]
{
  "code": "UNAUTHORIZED",
  "message": "Missing X-Tenant-ID header. Multi-tenant access requires a tenant identifier.",
  "correlationId": "",
  "timestamp": "2026-06-20T13:40:38.705363+00:00"
}
\end{lstlisting}

% ============================================================
\section{Cat\'alogo de errores y soluciones}
% ============================================================

\begin{tcolorbox}[colback=errred!5!white,colframe=errred,breakable,title={Error 1: iam:CreateRole denied}]
\textbf{Cuando}: \texttt{sam deploy} intenta crear el role de la Lambda. \\
\textbf{Causa}: AWS Academy no permite crear roles. \\
\textbf{Soluci\'on}: usar \texttt{LabRole} pre-existente (\texttt{Role: !Ref LambdaRoleArn}).
\end{tcolorbox}

\begin{tcolorbox}[colback=errred!5!white,colframe=errred,breakable,title={Error 2: iam:PassRole denied}]
\textbf{Cuando}: \texttt{sam deploy} con otro role que no sea \texttt{LabRole}. \\
\textbf{Causa}: el usuario voclabs solo puede pasar LabRole. \\
\textbf{Soluci\'on}: cambiar el ARN a \texttt{LabRole}.
\end{tcolorbox}

\begin{tcolorbox}[colback=errred!5!white,colframe=errred,breakable,title={Error 3: Tags denied en IAM resources}]
\textbf{Cuando}: el template tiene \texttt{Tags:} en una funci\'on. \\
\textbf{Causa}: Academy bloquea tagging en IAM. \\
\textbf{Soluci\'on}: quitar \texttt{Tags} del template.
\end{tcolorbox}

\begin{tcolorbox}[colback=errred!5!white,colframe=errred,breakable,title={Error 4: No module named 'src'}]
\textbf{Cuando}: handler con prefijo \texttt{src.handlers.*}. \\
\textbf{Causa}: el paquete \texttt{src} no est\'a en el path de importaci\'on. \\
\textbf{Soluci\'on}: handlers como top-level (\texttt{handlers.create\_queja.lambda\_handler}).
\end{tcolorbox}

\begin{tcolorbox}[colback=errred!5!white,colframe=errred,breakable,title={Error 5: email-validator not installed}]
\textbf{Cuando}: el body usa \texttt{EmailStr} de Pydantic. \\
\textbf{Causa}: Pydantic no incluye email-validator por defecto. \\
\textbf{Soluci\'on}: \texttt{pip install email-validator} y agregar a requirements.txt.
\end{tcolorbox}

\begin{tcolorbox}[colback=errred!5!white,colframe=errred,breakable,title={Error 6: KeyError 'message' in LogRecord}]
\textbf{Cuando}: \texttt{logger.warning(extra=\{'message': ...\})}. \\
\textbf{Causa}: \texttt{message} es atributo interno de logging. \\
\textbf{Soluci\'on}: usar \texttt{error\_message} u otro nombre.
\end{tcolorbox}

\begin{tcolorbox}[colback=errred!5!white,colframe=errred,breakable,title={Error 7: Stack en ROLLBACK\_COMPLETE}]
\textbf{Cuando}: primer deploy falla y no se puede re-deployar. \\
\textbf{Causa}: CloudFormation no permite update sobre ROLLBACK\_COMPLETE. \\
\textbf{Soluci\'on}:
\begin{lstlisting}
aws cloudformation delete-stack --stack-name sentinel-academia-dev
aws cloudformation wait stack-delete-complete --stack-name sentinel-academia-dev
sam deploy ...
\end{lstlisting}
\end{tcolorbox}

\begin{tcolorbox}[colback=errred!5!white,colframe=errred,breakable,title={Error 8: \_\_pycache\_\_ permission denied en rebuild}]
\textbf{Cuando}: \texttt{sam build} despu\'es de un build con Docker. \\
\textbf{Causa}: archivos creados por root en container Docker. \\
\textbf{Soluci\'on}: \texttt{sam build -{}-build-dir .aws-sam-new}
\end{tcolorbox}

% ============================================================
\section{Lo que aprendimos (resumen ejecutivo)}
% ============================================================

\begin{tcolorbox}[colback=okgreen!5!white,colframe=okgreen,title={\textbf{Checklist Fase 1}}]
\begin{enumerate}
  \item \textbf{1 Lambda = 1 archivo = 1 endpoint}
  \item \textbf{\_shared/} para c\'odigo com\'un (logger, errors, http, schemas)
  \item \textbf{LabRole} para todas las Lambdas (nunca crear roles)
  \item \textbf{CodeUri: ../src} apuntando a src/, no al project root
  \item \textbf{Handlers como top-level} (sin prefijo \texttt{src.})
  \item \textbf{202 Accepted} para operaciones as\'incronas
  \item \textbf{Error format spec-aligned}: \texttt{\{code, message, correlationId, timestamp\}}
  \item \textbf{Multi-tenant via X-Tenant-ID header} (nunca del body)
  \item \textbf{DynamoDB single-table} con PK \texttt{TENANT\#xxx\#QUEJA\#yyy}
  \item \textbf{email-validator} obligatorio para \texttt{EmailStr}
  \item \textbf{No usar 'message' en logger extra} (usar 'error\_message')
  \item \textbf{sam build -{}-build-dir .aws-sam-new} para evitar problemas de permisos
\end{enumerate}
\end{tcolorbox}

% ============================================================
\section{Ejercicios para practicar}
% ============================================================

\subsection{Ejercicio 1: A\~nadir un endpoint \texttt{GET /api/quejas/\{id\}/status}}

Endpoint que devuelve solo el status actual de una queja (no toda la data).

\textbf{Pista}: crea \filename{src/handlers/get\_queja\_status.py}, reutiliza \texttt{dynamo.get\_queja()}, y mapea el schema.

\subsection{Ejercicio 2: A\~nadir filtro por status en \texttt{list\_quejas}}

Modifica el handler para que acepte \texttt{?status=PENDIENTE} y use el GSI1.

\textbf{Pista}: cambia \texttt{query\_by\_tenant()} para que use el GSI cuando hay filtro de status.

\subsection{Ejercicio 3: A\~nadir paginaci\'on con cursor real}

Hoy \texttt{nextCursor: null}. Implementa cursor opaco basado en \texttt{LastEvaluatedKey} de DynamoDB.

\subsection{Ejercicio 4: Crear script de smoke tests}

Automatiza los 5 tests en \filename{tests/smoke.sh} con aserciones (jq) y exit codes.

\subsection{Ejercicio 5: Deploy con \texttt{sam sync}}

En vez de \texttt{sam deploy}, usa \texttt{sam sync -{}-stack-name ... -{}-watch} para iterar m\'as r\'apido.

% ============================================================
\section{Siguiente paso: Fase 2}
% ============================================================

En la \textbf{Fase 2} agregamos el LLM:

\begin{itemize}
  \item Nuevo servicio \texttt{src/services/llm\_client.py} que envuelve OCI Cohere
  \item Lambda \texttt{analyze\_queja.py} consumer de SQS
  \item SQS queue para procesamiento as\'incrono
  \item Funci\'on que toma una queja, llama Cohere, actualiza status a ANALIZADA
  \item El \texttt{createQueja} env\'ia mensaje a SQS en vez de procesar directo
\end{itemize}

\textbf{Tiempo estimado}: 3-4 horas.

\begin{center}
\begin{tabular}{ll}
\toprule
Fase & Estado \\
\midrule
0. Setup + OCI & \checkmark \\
1. Backend Core (este tutorial) & \checkmark \\
2. LLM integration & pr\'oximo \\
3. Frontend deploy & pendiente \\
4. Archivos + Email & pendiente \\
5. RAG & opcional \\
6. Testing & pendiente \\
7. Deploy final + demo & pendiente \\
\bottomrule
\end{tabular}
\end{center}

\vspace{2em}
\begin{center}
\textit{Fin del tutorial. ¡A por la Fase 2!}
\end{center}

\end{document}
